\documentclass[11pt,a4paper]{article}

% Packages
\usepackage[utf8]{inputenc}
\usepackage[margin=1in]{geometry}
\usepackage{graphicx}
\usepackage{amsmath,amssymb}
\usepackage{hyperref}
\usepackage{cite}
\usepackage{booktabs}
\usepackage{multirow}
\usepackage{xcolor}
\usepackage{algorithm}
\usepackage{algorithmic}
\usepackage{float}

% Hyperref setup
\hypersetup{
    colorlinks=true,
    linkcolor=blue,
    filecolor=magenta,      
    urlcolor=cyan,
    citecolor=blue,
}

% Title information
\title{\textbf{Vision-Language Models in Radiology:\\
Zero-Shot Classification, Cross-Modal Retrieval, and\\
Factual Radiology Report Generation}\\
\vspace{0.3cm}
\large Research Proposal for Biomedical Image Analysis and Processing}

\author{
Graduate Student\\
Department of Computer Engineering\\
Sharif University of Technology\\
\texttt{email@sharif.edu}
}

\date{January 2026}

\begin{document}

\maketitle

\begin{abstract}
Radiology report generation and interpretation are critical yet time-consuming tasks in clinical workflows. Recent advances in vision-language models (VLMs), particularly contrastive learning approaches like CLIP, have demonstrated promising capabilities for zero-shot medical image classification and cross-modal retrieval. However, significant challenges remain: (1) factual hallucination in generated reports, where models produce clinically incorrect findings; (2) limited cross-institutional robustness, with performance degradation on external datasets; and (3) lack of uncertainty quantification for clinical deployment. This research proposes a comprehensive study of VLMs for radiology, focusing on three interconnected tasks: zero-shot pathology classification, image-text retrieval, and factuality-aware report generation. Building on CheXzero and R2Gen as foundational architectures, we propose novel extensions incorporating uncertainty quantification, RadGraph-based factuality constraints, and systematic cross-dataset evaluation. Using publicly available datasets (IU X-Ray, NIH ChestX-ray14), we aim to develop methods that advance both technical performance and clinical deployability. This work is designed as a graduate course project with clear pathways to publication at MICCAI, MIDL, or related venues, addressing critical gaps in medical AI research.
\end{abstract}

\tableofcontents
\newpage

%=====================================================
\section{Introduction}
%=====================================================

\subsection{Clinical Motivation}

Radiology is the cornerstone of modern medical diagnosis, with over 3 billion imaging studies performed annually worldwide. Each chest X-ray examination requires expert interpretation to detect pathologies ranging from common conditions (pneumonia, pleural effusion) to life-threatening emergencies (pneumothorax, aortic dissection). However, the radiology workforce faces unprecedented challenges:

\begin{itemize}
    \item \textbf{Global radiologist shortage:} The World Health Organization estimates a deficit of over 2 million radiologists globally, particularly severe in low- and middle-income countries.
    \item \textbf{Increasing imaging volumes:} Annual growth rate of 5-10\% in imaging studies outpaces radiologist training capacity.
    \item \textbf{Diagnostic errors:} Studies report 3-5\% error rates in radiology interpretation, with significant consequences for patient outcomes.
    \item \textbf{Workflow inefficiencies:} Radiologists spend 60-70\% of their time on report writing rather than complex diagnostic reasoning.
\end{itemize}

Artificial intelligence, specifically deep learning for medical image analysis, has emerged as a promising solution. Early successes in classification tasks (e.g., CheXNet achieving radiologist-level pneumonia detection) demonstrated technical feasibility. However, clinical translation remains limited due to fundamental challenges in trustworthiness, generalizability, and practical deployment.

\subsection{The Promise of Vision-Language Models}

Recent advances in vision-language models (VLMs) offer a paradigm shift from traditional supervised learning. Instead of requiring thousands of manually labeled examples per disease, VLMs learn joint representations of images and text through self-supervised contrastive learning on naturally occurring image-report pairs. This approach provides several advantages:

\begin{enumerate}
    \item \textbf{Zero-shot learning:} Models can classify new diseases without disease-specific training data, using only text descriptions.
    \item \textbf{Data efficiency:} Self-supervised pretraining on unlabeled image-text pairs reduces annotation burden.
    \item \textbf{Unified framework:} Single model supports multiple tasks (classification, retrieval, generation).
    \item \textbf{Interpretability:} Text grounding provides clinically meaningful explanations.
\end{enumerate}

Models like CLIP (Contrastive Language-Image Pre-training) have revolutionized natural image understanding. Medical adaptations—CheXzero, GLORIA, BioViL—demonstrate that similar principles apply to radiology, achieving expert-level performance on pathology detection tasks.

\subsection{Critical Research Gaps}

Despite promising results, current medical VLMs face severe limitations that impede clinical deployment:

\textbf{Gap 1: Factual Hallucination}\\
VLMs trained on report generation frequently produce fluent but clinically incorrect text. A model might generate "fracture of the left 4th rib" when no fracture exists, or omit critical findings like pneumothorax. Traditional NLG metrics (BLEU, ROUGE) fail to detect these errors, as they measure linguistic similarity, not clinical accuracy.

\textbf{Gap 2: Cross-Institutional Brittleness}\\
Models trained on single-institution data (e.g., MIMIC-CXR from Beth Israel Deaconess) show performance degradation of 5-15\% AUROC when tested on external hospitals. Contributing factors include equipment differences, population demographics, and reporting style variations.

\textbf{Gap 3: Lack of Uncertainty Quantification}\\
Clinical decisions require knowing when a model is uncertain. Current VLMs produce point predictions without confidence estimates, making it impossible to identify out-of-distribution cases or defer to human experts appropriately.

\textbf{Gap 4: Limited Grounding}\\
Many models fail to link generated findings to image regions. Without visual evidence for each claim, radiologists cannot verify model outputs or understand failure modes.

\textbf{Gap 5: Evaluation Misalignment}\\
Standard metrics (BLEU, ROUGE) correlate poorly (r=0.3) with clinical utility. Emerging alternatives like RadGraph F1 and CheXbert accuracy better capture clinical correctness but remain underutilized in training objectives.

\subsection{Research Objectives}

This project aims to address these gaps through three primary objectives:

\textbf{Objective 1: Develop uncertainty-aware zero-shot classification}\\
Extend CheXzero with Monte Carlo Dropout and conformal prediction to provide calibrated confidence intervals, enabling appropriate human-AI collaboration.

\textbf{Objective 2: Create factuality-constrained report generation}\\
Integrate RadGraph-based supervision into R2Gen training, penalizing entity hallucination and rewarding grounded generation.

\textbf{Objective 3: Establish cross-dataset robustness benchmarks}\\
Systematically evaluate model performance across multiple public datasets (IU X-Ray, NIH ChestX-ray14), analyze failure modes, and propose domain adaptation techniques.

\subsection{Scope and Constraints}

This project focuses on \textbf{publicly available resources only}:
\begin{itemize}
    \item \textbf{Datasets:} IU X-Ray (3,955 studies), NIH ChestX-ray14 (112,120 images)
    \item \textbf{Modality:} Chest X-ray (frontal view)
    \item \textbf{Tasks:} Classification, retrieval, report generation
    \item \textbf{Compute:} Single GPU experiments (NVIDIA V100 or A100)
\end{itemize}

While MIMIC-CXR (227K studies) offers larger scale, its credentialing requirements and 1-2 week access delay make it impractical for an 8-month course project. Our approach demonstrates that meaningful research contributions are possible with smaller, fully public datasets by focusing on methodological innovation rather than scale.

\subsection{Expected Contributions}

This research will deliver:

\begin{enumerate}
    \item \textbf{Technical contributions:}
    \begin{itemize}
        \item Uncertainty-aware VLM architecture with calibrated predictions
        \item RadGraph-guided training loss for factual report generation
        \item Cross-dataset evaluation protocol and benchmark results
    \end{itemize}
    
    \item \textbf{Scientific contributions:}
    \begin{itemize}
        \item Systematic analysis of VLM robustness across datasets
        \item Quantification of factuality improvement via knowledge graph supervision
        \item Ablation studies isolating key design choices
    \end{itemize}
    
    \item \textbf{Practical contributions:}
    \begin{itemize}
        \item Open-source implementation with pretrained weights
        \item Reproducible evaluation framework
        \item Clinical error analysis and failure mode taxonomy
    \end{itemize}
    
    \item \textbf{Publication:}
    \begin{itemize}
        \item Target venue: MICCAI, MIDL, ML4H Workshop, or TMI journal
        \item Estimated timeline: Submission by Month 7-8
    \end{itemize}
\end{enumerate}

\subsection{Document Organization}

The remainder of this proposal is organized as follows: Section 2 reviews related work in medical VLMs, report generation, and evaluation methods. Section 3 describes datasets and their characteristics. Section 4 presents our proposed methodology with three research extensions. Section 5 details the experimental plan and evaluation protocol. Section 6 discusses expected contributions and limitations. Section 7 addresses ethical considerations. Section 8 provides a detailed timeline. Section 9 concludes with future directions.

%=====================================================
\section{Related Work}
%=====================================================

\subsection{Contrastive Vision-Language Learning}

\subsubsection{Natural Image VLMs}

The CLIP model \cite{radford2021learning} pioneered large-scale contrastive learning between images and text, demonstrating remarkable zero-shot transfer to downstream tasks. Trained on 400M image-text pairs from the internet, CLIP learns to align visual and language representations in a shared embedding space. The key insight: maximizing agreement between positive pairs (matching image-text) while minimizing agreement with negative pairs (non-matching) creates semantically meaningful representations.

CLIP's architecture consists of:
\begin{itemize}
    \item \textbf{Image Encoder:} Vision Transformer (ViT) or ResNet
    \item \textbf{Text Encoder:} Transformer (GPT-style)
    \item \textbf{Training Objective:} Symmetric cross-entropy over cosine similarities
\end{itemize}

Zero-shot classification emerges naturally: given image $I$ and candidate labels $\{L_1, ..., L_N\}$, embed each as prompt "a photo of [label]", compute $\text{sim}(I, L_i)$, and select $\arg\max_i \text{sim}(I, L_i)$.

However, direct application of CLIP to medical imaging fails due to domain mismatch: (1) medical images differ fundamentally from natural photos, (2) clinical terminology absent from internet text, and (3) diagnostic reasoning requires fine-grained discrimination beyond object recognition.

\subsubsection{Medical VLM Adaptations}

\textbf{ConVIRT} \cite{zhang2020convirt} was among the first to adapt contrastive learning to radiology. Pretrained on MIMIC-CXR's 377K image-report pairs, ConVIRT achieves +3.2\% AUROC improvement over ImageNet initialization on CheXpert classification. Key findings: (1) medical-specific pretraining outperforms natural image pretraining, (2) gains largest in few-shot regimes (1-10\% labeled data), (3) transfer learning effective across chest X-ray datasets.

\textbf{CheXzero} \cite{tiu2022chexzero} demonstrates expert-level zero-shot pathology detection. Using ResNet-50 + BioClinicalBERT, CheXzero achieves 0.873 mean AUROC on 5 CheXpert diseases without disease-specific labels. Critical innovation: carefully engineered clinical prompts ("Findings consistent with pneumonia" vs. "No evidence of pneumonia"). Limitations: prompt engineering requires clinical expertise and manual trial-and-error.

\textbf{GLORIA} \cite{huang2021gloria} extends global contrastive learning with local attention, matching image regions to report sentences. Achieves state-of-the-art retrieval (32.2\% Recall@1 on MIMIC-CXR) and provides interpretable attention maps. Local alignment enables grounding but increases computational cost (4× V100 GPUs, 2-3 days training).

\textbf{BioViL} \cite{microsoft2022biovil} (Microsoft, ECCV 2022) uses RadGraph to create structured image-text alignments during pretraining. By parsing reports into knowledge graphs and aligning entities to image regions, BioViL improves phrase grounding accuracy by 20\% over baseline. Demonstrates value of structured medical knowledge in VLM training.

\textbf{MedCLIP, PubMedCLIP, PMC-CLIP} explore pretraining on broader medical data (PubMed figures, captions) for multi-domain transfer. While achieving reasonable zero-shot performance across modalities (CT, MRI, microscopy), they lack radiology-specific optimization of specialized models like CheXzero.

\subsection{Radiology Report Generation}

Automatic report generation aims to produce coherent, clinically accurate radiology reports from images. Unlike image captioning, medical reports require:
\begin{itemize}
    \item Structured format (Findings, Impressions, Comparisons)
    \item Clinical precision (diagnostic terminology, negation handling)
    \item Evidence linking (findings supported by visible pathology)
    \item Absence documentation ("no evidence of X" equally important)
\end{itemize}

\subsubsection{Encoder-Decoder Architectures}

Early approaches adapted Show-Attend-Tell \cite{xu2015show} from image captioning. \textbf{TieNet} \cite{wang2018tienet} uses multi-task learning (generation + classification) to improve clinical accuracy. \textbf{AdaAtt} applies adaptive attention to focus on relevant regions. These methods achieve BLEU-4 $\approx$ 0.18-0.20 on IU X-Ray but suffer from generic descriptions and factual errors.

\textbf{R2Gen} \cite{chen2020r2gen} introduces a relational memory module storing clinical patterns. Memory slots specialize to common findings ("lungs clear", "cardiomegaly"), improving coherence. Achieves BLEU-4 = 0.203 on IU X-Ray, CIDEr = 0.485. Multi-scale visual encoding (CNN + Transformer) captures both global and local features. Despite strong linguistic metrics, R2Gen lacks explicit factuality constraints.

\textbf{CMN (Clinical Memory Network)} \cite{chen2021cmn} extends memory-based generation with clinical prototypes retrieved from training data. Achieves modest CheXbert F1 improvement but limited evaluation of factual correctness.

\textbf{PPKED} \cite{liu2022ppked} integrates medical knowledge graphs (RadLex, UMLS) to guide generation. While conceptually appealing, complex pipeline and limited code availability hinder reproducibility.

\subsubsection{Retrieval-Based and Hybrid Methods}

\textbf{RATCHET} \cite{hou2022ratchet} retrieves similar historical reports and adapts them to current image. Achieves high fluency but limited novelty—struggles with rare findings absent from retrieval database.

\textbf{Recent trends (2023-2024):} Large language model (LLM) integration. LLaVA-Med, Med-Flamingo, and RadFM use instruction-tuned LLMs (LLaMA, GPT-3.5) conditioned on visual features. While impressive in qualitative examples, quantitative factuality evaluation remains limited.

\subsection{Factuality and Grounding}

\subsubsection{Evaluation Beyond BLEU/ROUGE}

Traditional NLG metrics (BLEU, ROUGE, METEOR, CIDEr) measure n-gram overlap with reference text. Critical flaw: high scores possible with clinically incorrect content. Example: "mild cardiomegaly" vs. "severe cardiomegaly" have high BLEU (8/9 words match) but opposite clinical meaning.

\textbf{RadGraph} \cite{jain2023radgraph} extracts clinical entities and relations from reports, converting text to knowledge graphs. RadGraph F1 metric compares hypothesis and reference graphs, capturing clinical content independently of phrasing. Correlation with expert assessment: r=0.7 (vs. BLEU r=0.3). Now standard evaluation for medical report generation.

\textbf{CheXbert} \cite{smit2020chexbert} fine-tunes BERT to classify 14 CheXpert labels from report text. CheXbert accuracy measures whether generated reports convey correct diagnostic findings. Simpler than RadGraph but limited to 14 binary labels.

\subsubsection{Grounding and Attention}

\textbf{MS-CXR} \cite{boecking2022ms} provides sentence-to-region annotations, enabling supervised grounding. Models learn to link "opacity in left lower lobe" to specific image coordinates.

\textbf{GREEN} \cite{liu2022green} uses entity-level supervision to reduce hallucination. By explicitly annotating medical entities (anatomy, observations) and penalizing ungrounded generation, GREEN achieves 15\% reduction in false entity rate.

\textbf{FactualGAN} \cite{zhou2022factual} trains adversarial discriminator to distinguish factual from hallucinated reports, providing training signal for generator.

\subsection{Cross-Dataset Evaluation and Domain Shift}

Most medical VLM papers evaluate on a single dataset (typically MIMIC-CXR or IU X-Ray). Limited cross-dataset validation obscures generalization failure. Known issues:

\begin{itemize}
    \item \textbf{Equipment variation:} GE vs. Siemens scanners produce different image characteristics
    \item \textbf{Population demographics:} Age, comorbidity distributions vary across institutions
    \item \textbf{Reporting style:} Phrase choices, sentence structure, abbreviations differ
    \item \textbf{Disease prevalence:} COVID-19, endemic diseases, referral patterns shift distributions
\end{itemize}

Few studies systematically evaluate robustness:
\begin{itemize}
    \item CheXzero tests on PadChest (Spain) and observes 3-5\% AUROC drop
    \item GLORIA shows modest degradation on NIH-14
    \item Most papers ignore external validation entirely
\end{itemize}

\textbf{Research gap:} Comprehensive cross-dataset benchmarking with analysis of failure modes and domain adaptation techniques.

\subsection{Uncertainty Quantification in Medical AI}

Clinical deployment requires knowing \textit{when} a model is uncertain. Uncertainty quantification (UQ) methods:

\begin{itemize}
    \item \textbf{Bayesian Deep Learning:} Monte Carlo Dropout \cite{gal2016dropout}, variational inference
    \item \textbf{Ensemble Methods:} Train multiple models, report disagreement
    \item \textbf{Conformal Prediction:} Distribution-free confidence sets with coverage guarantees
    \item \textbf{Calibration:} Temperature scaling, Platt scaling for probability calibration
\end{itemize}

Medical imaging applications:
\begin{itemize}
    \item Dermoscopy: Uncertainty-aware skin lesion classification defers 15\% of ambiguous cases to dermatologists, improving overall accuracy
    \item Pathology: Conformal prediction provides slide-level confidence intervals for tumor classification
    \item Radiology: Limited work on VLM uncertainty—most models produce overconfident predictions
\end{itemize}

\textbf{Research gap:} Integration of principled UQ into medical VLMs for safe clinical deployment.

\subsection{Summary and Positioning}

Existing work demonstrates:
\begin{itemize}
    \item Contrastive VLMs achieve strong zero-shot classification (CheXzero: 0.87 AUROC)
    \item Report generation reaches moderate linguistic quality (R2Gen: BLEU-4 = 0.20)
    \item Clinical evaluation metrics (RadGraph, CheXbert) quantify factual correctness
\end{itemize}

Critical gaps remain:
\begin{itemize}
    \item Factual hallucination unaddressed in training objectives
    \item Cross-dataset robustness poorly understood
    \item Uncertainty quantification absent
    \item Grounding between generated text and image regions weak
\end{itemize}

\textbf{Our contribution:} This project bridges these gaps by:
\begin{enumerate}
    \item Incorporating RadGraph-based factuality loss in report generation training
    \item Systematic cross-dataset evaluation (IU X-Ray $\leftrightarrow$ NIH-14)
    \item Uncertainty quantification via MC Dropout and conformal prediction
    \item Public dataset focus, enabling reproducible research
\end{enumerate}

Next section describes datasets in detail.

%=====================================================
\section{Datasets}
%=====================================================

This section provides comprehensive descriptions of publicly available datasets used in our research. All datasets require no institutional access or credentialing beyond simple registration (if any).

\subsection{IU X-Ray (Indiana University Chest X-Ray Collection)}

\subsubsection{Overview}
\begin{itemize}
    \item \textbf{Institution:} Indiana University School of Medicine
    \item \textbf{Year:} 2015
    \item \textbf{Access:} Fully public via OpenI \url{https://openi.nlm.nih.gov/}
    \item \textbf{Size:} 7,470 images, 3,955 studies, 3,955 patients
    \item \textbf{File Size:} $\sim$7 GB images + 5 MB reports (XML)
\end{itemize}

\subsubsection{Data Composition}

\textbf{Images:}
\begin{itemize}
    \item Modality: Chest X-ray (frontal PA/AP + lateral)
    \item 40\% frontal-only, 60\% frontal+lateral
    \item Average: 1.89 images per study
    \item Format: PNG, grayscale, resolution 2048$\times$2048 to 3000$\times$2500
\end{itemize}

\textbf{Reports:} Each study has structured radiology report with 4 sections:
\begin{enumerate}
    \item \textbf{Comparison:} Reference to prior studies (if available)
    \item \textbf{Indication:} Clinical reason for examination
    \item \textbf{Findings:} Detailed observations of chest anatomy and pathology
    \item \textbf{Impression:} Summary diagnosis and recommendations
\end{enumerate}

Average report length: 35-40 words. Characteristics: relatively concise, high frequency of negation ("no acute disease"), some PHI redacted with placeholders.

\textbf{Annotations:}
\begin{itemize}
    \item Manual expert labels for 14 observations: Normal (28.1\%), Cardiomegaly (9.5\%), Lung Opacity (16.6\%), Edema (2.2\%), Consolidation (4.4\%), Pneumonia (4.0\%), Atelectasis (5.1\%), Pneumothorax (1.1\%), Pleural Effusion (7.7\%), etc.
    \item Severity grades: Mild, Moderate, Severe
    \item MeSH (Medical Subject Headings) tags: Controlled vocabulary for retrieval
\end{itemize}

\subsubsection{Data Splits}
Standard split used in literature (e.g., R2Gen, CMN):
\begin{center}
\begin{tabular}{lcc}
\toprule
\textbf{Split} & \textbf{Studies} & \textbf{Percentage} \\
\midrule
Train & 2,770 & 70\% \\
Validation & 585 & 15\% \\
Test & 600 & 15\% \\
\bottomrule
\end{tabular}
\end{center}

Random split at study level. We adopt this standard for comparability.

\subsubsection{Strengths and Limitations}

\textbf{Strengths:}
\begin{itemize}
    \item Only large-scale public dataset with paired images and reports
    \item Well-structured reports enable easy parsing
    \item Manual disease annotations provide gold-standard labels
    \item Standard benchmark for report generation ($>$100 papers)
    \item Multi-view data (frontal + lateral) provides richer information
\end{itemize}

\textbf{Limitations:}
\begin{itemize}
    \item Small scale (3,955 vs. MIMIC-CXR's 227K) limits model capacity
    \item Single institution $\rightarrow$ potential bias in equipment, population, reporting style
    \item No demographic data (age, gender, race) $\rightarrow$ cannot study fairness
    \item Short reports (35 words) vs. typical clinical reports (50-100 words)
    \item Class imbalance: rare diseases like pneumothorax (1.1\%) underrepresented
\end{itemize}

\subsubsection{Use in This Project}

\textbf{Primary dataset for:}
\begin{itemize}
    \item Training vision-language models (contrastive pretraining)
    \item Report generation experiments (R2Gen baseline and extensions)
    \item Image-text retrieval evaluation
\end{itemize}

\textbf{Preprocessing:}
\begin{enumerate}
    \item Parse XML to extract findings and impression sections
    \item Resize images to 224$\times$224 (CheXzero) or 320$\times$320
    \item Normalize with ImageNet statistics or dataset-specific mean/std
    \item Tokenize text with BioClinicalBERT tokenizer (max 128 tokens)
\end{enumerate}

\subsection{NIH ChestX-ray14}

\subsubsection{Overview}
\begin{itemize}
    \item \textbf{Institution:} National Institutes of Health Clinical Center
    \item \textbf{Year:} 2017
    \item \textbf{Access:} Fully public \url{https://nihcc.app.box.com/v/ChestXray-NIHCC}
    \item \textbf{Size:} 112,120 images from 30,805 patients
    \item \textbf{File Size:} $\sim$45 GB (uncompressed PNG)
\end{itemize}

\subsubsection{Data Composition}

\textbf{Images:}
\begin{itemize}
    \item Modality: Chest X-ray (frontal view only: PA/AP)
    \item Format: PNG, 1024$\times$1024, 8-bit grayscale
    \item Standardized resolution (unlike IU X-Ray's variable sizes)
\end{itemize}

\textbf{Labels:} 14 disease categories extracted via NLP from original radiology reports (reports not publicly released):

\begin{center}
\small
\begin{tabular}{lcc}
\toprule
\textbf{Disease} & \textbf{Positive Cases} & \textbf{Prevalence} \\
\midrule
Infiltration & 19,894 & 17.7\% \\
Effusion & 13,317 & 11.9\% \\
Atelectasis & 11,559 & 10.3\% \\
Nodule & 6,331 & 5.6\% \\
Pneumothorax & 5,302 & 4.7\% \\
Mass & 5,782 & 5.2\% \\
Consolidation & 4,667 & 4.2\% \\
Pleural Thickening & 3,385 & 3.0\% \\
Cardiomegaly & 2,776 & 2.5\% \\
Emphysema & 2,516 & 2.2\% \\
Edema & 2,303 & 2.1\% \\
Fibrosis & 1,686 & 1.5\% \\
Pneumonia & 1,431 & 1.3\% \\
Hernia & 227 & 0.2\% \\
\textbf{No Finding} & 60,361 & 53.8\% \\
\bottomrule
\end{tabular}
\end{center}

\textbf{Important:} Labels are \textit{automatically extracted via NLP}, not manually annotated. Estimated 10-20\% label noise due to negation errors, context misunderstanding, and mention vs. diagnosis ambiguity.

\textbf{Bounding Boxes:} Subset of 880 images have radiologist-drawn bounding boxes for 8 diseases (Atelectasis, Cardiomegaly, Effusion, Infiltration, Mass, Nodule, Pneumonia, Pneumothorax). Enables weakly supervised localization.

\subsubsection{Data Splits}

Common split:
\begin{center}
\begin{tabular}{lcc}
\toprule
\textbf{Split} & \textbf{Images} & \textbf{Percentage} \\
\midrule
Train & 70,000 & 62\% \\
Validation & 17,120 & 15\% \\
Test & 25,000 & 23\% \\
\bottomrule
\end{tabular}
\end{center}

Note: Random split at image level (not patient level). Some papers use patient-level splits to avoid data leakage.

\subsubsection{Strengths and Limitations}

\textbf{Strengths:}
\begin{itemize}
    \item Large scale (112K images) enables robust deep learning
    \item Fully public with no access barriers
    \item Standardized 1024$\times$1024 resolution simplifies preprocessing
    \item Widely used benchmark ($>$500 citing papers)
    \item Bounding box subset enables localization research
    \item Diverse patient population (NIH Clinical Center)
\end{itemize}

\textbf{Limitations:}
\begin{itemize}
    \item \textbf{No reports:} Cannot use for report generation or VLM pretraining
    \item Label noise (10-20\% error rate) due to NLP extraction
    \item Extreme class imbalance (Hernia: 0.2\% vs. No Finding: 53.8\%)
    \item Frontal view only (no lateral)
    \item Single institution bias
    \item No demographic data in public release
\end{itemize}

\subsubsection{Use in This Project}

\textbf{Primary use cases:}
\begin{itemize}
    \item \textbf{Zero-shot classification evaluation:} Test VLMs pretrained on IU X-Ray without fine-tuning on NIH-14. Measures cross-dataset generalization.
    \item \textbf{Cross-dataset robustness:} Train on IU X-Ray (with reports), test on NIH-14 (classification only). Quantify performance drop.
    \item \textbf{Weakly supervised localization:} Evaluate attention maps on bounding box subset.
\end{itemize}

\textbf{Not suitable for:}
\begin{itemize}
    \item Report generation (no reports available)
    \item Image-text retrieval (no text beyond 14 labels)
\end{itemize}

\subsection{Dataset Comparison and Selection Rationale}

Table \ref{tab:dataset_comparison} compares key characteristics:

\begin{table}[h]
\centering
\caption{Comparison of Public Chest X-Ray Datasets}
\label{tab:dataset_comparison}
\small
\begin{tabular}{lccccc}
\toprule
\textbf{Dataset} & \textbf{Images} & \textbf{Reports} & \textbf{Labels} & \textbf{Access} & \textbf{Best For} \\
\midrule
IU X-Ray & 7.5K & \checkmark & Manual & Public & Generation \\
NIH-14 & 112K & \texttimes & NLP & Public & Classification \\
MIMIC-CXR & 377K & \checkmark & NLP & Credentialed & Pretraining \\
CheXpert & 224K & \texttimes & NLP & Registration & Multi-label \\
PadChest & 160K & Spanish & 174 labels & Public & European \\
\bottomrule
\end{tabular}
\end{table}

\textbf{Why IU X-Ray + NIH-14?}
\begin{enumerate}
    \item \textbf{Complementary:} IU X-Ray provides reports for generation; NIH-14 provides scale for classification
    \item \textbf{Fully public:} No credentialing delays or usage restrictions
    \item \textbf{Standard benchmarks:} Both widely used, enabling comparison with prior work
    \item \textbf{Sufficient for graduate project:} Combined 115K images adequate for meaningful experiments
\end{enumerate}

\textbf{Limitations accepted:}
\begin{itemize}
    \item Cannot match performance of models trained on MIMIC-CXR (10$\times$ larger)
    \item Must frame contributions as methodological innovations rather than state-of-the-art performance
    \item Focus on robustness, factuality, and uncertainty rather than raw accuracy
\end{itemize}

\subsection{Ethical Considerations for Datasets}

Both datasets de-identified per HIPAA standards:
\begin{itemize}
    \item Patient names, dates, IDs removed
    \item Ages >89 redacted to "90+"
    \item Geocodes broader than state level removed
\end{itemize}

Institutional Review Board (IRB) approval obtained by original dataset creators. Our use constitutes secondary analysis of public data, typically exempt from additional IRB review.

\textbf{Potential biases:}
\begin{itemize}
    \item Single-institution datasets may not generalize to global populations
    \item Referral patterns at academic medical centers differ from community hospitals
    \item Unknown demographic distributions limit fairness analysis
\end{itemize}

We acknowledge these limitations and will discuss generalization constraints in our final paper.

%=====================================================
\section{Proposed Methodology}
%=====================================================

\subsection{Overview}

Our approach combines two foundational architectures:

\begin{enumerate}
    \item \textbf{CheXzero} \cite{tiu2022chexzero}: Contrastive vision-language model for zero-shot classification and retrieval
    \item \textbf{R2Gen} \cite{chen2020r2gen}: Memory-driven transformer for report generation
\end{enumerate}

We propose three novel extensions addressing critical gaps:

\begin{itemize}
    \item \textbf{Extension 1 (Classification):} Uncertainty-aware zero-shot classification via Monte Carlo Dropout and conformal prediction
    \item \textbf{Extension 2 (Generation):} Factuality-constrained report generation using RadGraph-based training loss
    \item \textbf{Extension 3 (Robustness):} Cross-dataset evaluation protocol with systematic failure analysis
\end{itemize}

Each extension is designed to be:
\begin{itemize}
    \item \textbf{Feasible:} Implementable within 2-3 months on single GPU
    \item \textbf{Novel:} Addresses underexplored research gap
    \item \textbf{Publishable:} Sufficient contribution for MICCAI/MIDL workshop or conference paper
\end{itemize}

\subsection{Baseline: CheXzero for Zero-Shot Classification}

\subsubsection{Architecture}

CheXzero consists of:

\textbf{Image Encoder $f_v$:}
\begin{itemize}
    \item ResNet-50 pretrained on ImageNet
    \item Input: $x \in \mathbb{R}^{3 \times 320 \times 320}$ (chest X-ray, replicated to 3 channels)
    \item Output: $v = f_v(x) \in \mathbb{R}^{2048}$ (global average pooled features)
\end{itemize}

\textbf{Text Encoder $f_u$:}
\begin{itemize}
    \item BioClinicalBERT (BERT pretrained on PubMed + MIMIC notes)
    \item Input: $t$ (radiology report text, max 512 tokens)
    \item Output: $u = f_u(t) \in \mathbb{R}^{768}$ ([CLS] token embedding)
\end{itemize}

\textbf{Projection Heads:}
\begin{itemize}
    \item Linear layers project to shared 512-dim embedding space
    \item $z_v = W_v v$, $z_u = W_u u$, where $z_v, z_u \in \mathbb{R}^{512}$
    \item L2 normalization: $\tilde{z}_v = z_v / \|z_v\|_2$
\end{itemize}

\subsubsection{Training: Contrastive Learning}

Given batch of $N$ image-report pairs $\{(x_i, t_i)\}_{i=1}^N$:

\textbf{Contrastive Loss (Symmetric):}
\begin{equation}
\mathcal{L}_{\text{CL}} = -\frac{1}{2N} \sum_{i=1}^{N} \left[ \log \frac{\exp(\tilde{z}_{v,i}^T \tilde{z}_{u,i} / \tau)}{\sum_{j=1}^{N} \exp(\tilde{z}_{v,i}^T \tilde{z}_{u,j} / \tau)} + \log \frac{\exp(\tilde{z}_{u,i}^T \tilde{z}_{v,i} / \tau)}{\sum_{j=1}^{N} \exp(\tilde{z}_{u,j}^T \tilde{z}_{v,i} / \tau)} \right]
\end{equation}

where $\tau = 0.07$ is temperature parameter.

\textbf{Training Details:}
\begin{itemize}
    \item Optimizer: Adam, lr = $5 \times 10^{-5}$
    \item Batch size: 128 (or smaller with gradient accumulation)
    \item Epochs: 15-20
    \item Data augmentation: Random horizontal flip, rotation $\pm 10^\circ$
    \item Hardware: 1$\times$ NVIDIA A100 (24 hours) or V100 (48 hours)
\end{itemize}

\subsubsection{Zero-Shot Classification}

For disease $d$, construct prompt templates:

\textbf{Positive prompts $\mathcal{T}_+$:}
\begin{itemize}
    \item "Findings consistent with [disease]"
    \item "Evidence of [disease]"
    \item "Patient has [disease]"
\end{itemize}

\textbf{Negative prompts $\mathcal{T}_-$:}
\begin{itemize}
    \item "No evidence of [disease]"
    \item "Findings not consistent with [disease]"
    \item "Normal chest X-ray without [disease]"
\end{itemize}

\textbf{Classification score:}
\begin{equation}
s_d(x) = \frac{1}{|\mathcal{T}_+|} \sum_{t \in \mathcal{T}_+} \cos(\tilde{z}_v(x), \tilde{z}_u(t)) - \frac{1}{|\mathcal{T}_-|} \sum_{t \in \mathcal{T}_-} \cos(\tilde{z}_v(x), \tilde{z}_u(t))
\end{equation}

Decision: $\hat{y}_d = \mathbb{1}[s_d(x) > \theta]$, where $\theta$ is threshold (typically 0, or calibrated on validation set).

\subsubsection{Expected Baseline Performance}

Based on CheXzero paper and our preliminary experiments:

\begin{center}
\begin{tabular}{lcc}
\toprule
\textbf{Disease} & \textbf{IU X-Ray (AUROC)} & \textbf{NIH-14 (AUROC)} \\
\midrule
Atelectasis & 0.78-0.81 & 0.73-0.76 \\
Cardiomegaly & 0.83-0.86 & 0.80-0.83 \\
Edema & 0.88-0.91 & 0.85-0.88 \\
Pleural Effusion & 0.90-0.93 & 0.87-0.90 \\
\textbf{Mean (5 diseases)} & \textbf{0.84-0.87} & \textbf{0.81-0.84} \\
\bottomrule
\end{tabular}
\end{center}

Performance drop on NIH-14 expected due to domain shift (single institution vs. NIH Clinical Center).

\subsection{Extension 1: Uncertainty-Aware Classification}

\subsubsection{Motivation}

Clinical deployment requires knowing when model is uncertain. Current CheXzero produces deterministic scores without confidence estimates. Consequences:
\begin{itemize}
    \item Cannot defer ambiguous cases to human experts
    \item Overconfident predictions on out-of-distribution data
    \item Difficult to set operating points balancing sensitivity and specificity
\end{itemize}

\textbf{Research Question:} Can we augment CheXzero with calibrated uncertainty estimates while maintaining zero-shot performance?

\subsubsection{Technical Approach}

We propose two complementary methods:

\textbf{Method 1: Monte Carlo Dropout (MC Dropout)}

\textbf{Idea:} Treat dropout as Bayesian approximation. At inference, enable dropout and perform $T$ forward passes:

\begin{equation}
\hat{p}(y_d=1|x) = \frac{1}{T} \sum_{t=1}^{T} \sigma(s_d^{(t)}(x))
\end{equation}

where $s_d^{(t)}(x)$ is classification score with dropout mask $t$, $\sigma$ is sigmoid.

\textbf{Predictive Entropy (Uncertainty):}
\begin{equation}
H(y_d|x) = -\hat{p} \log \hat{p} - (1-\hat{p}) \log (1-\hat{p})
\end{equation}

High entropy $\rightarrow$ model uncertain $\rightarrow$ defer to radiologist.

\textbf{Implementation:}
\begin{itemize}
    \item Add dropout (p=0.2) after penultimate layer of ResNet-50 and BERT
    \item At inference: $T=20$ stochastic forward passes (overhead: 20$\times$ single pass)
    \item Aggregate predictions: mean $\pm$ std
\end{itemize}

\textbf{Method 2: Conformal Prediction}

\textbf{Idea:} Distribution-free confidence sets with coverage guarantees. Given calibration set $\mathcal{D}_{\text{cal}} = \{(x_i, y_i)\}$, construct prediction sets $C(x)$ such that:
\begin{equation}
\mathbb{P}(y \in C(x)) \geq 1 - \alpha
\end{equation}
for user-specified error rate $\alpha$ (e.g., 0.05 for 95\% coverage).

\textbf{Procedure:}
\begin{enumerate}
    \item Compute non-conformity scores on calibration set: $R_i = |s_d(x_i) - y_i|$
    \item Compute quantile: $q = \text{Quantile}(R_1, ..., R_m; 1-\alpha)$
    \item For test image $x$: predict $\hat{y} \in [s_d(x) - q, s_d(x) + q]$
\end{enumerate}

\textbf{Advantages:}
\begin{itemize}
    \item Finite-sample coverage guarantee (not asymptotic)
    \item No distributional assumptions
    \item Fast calibration (single pass over calibration set)
\end{itemize}

\textbf{Disadvantage:} Requires held-out calibration set (use validation split).

\subsubsection{Evaluation}

\textbf{Metrics:}
\begin{enumerate}
    \item \textbf{Calibration:} Expected Calibration Error (ECE)
    \begin{equation}
    \text{ECE} = \sum_{m=1}^{M} \frac{|B_m|}{N} |\text{acc}(B_m) - \text{conf}(B_m)|
    \end{equation}
    where $B_m$ are bins of predicted confidence, $\text{acc}(B_m)$ is accuracy in bin, $\text{conf}(B_m)$ is average confidence.
    
    \item \textbf{Coverage:} For conformal prediction, verify empirical coverage $\geq 1-\alpha$.
    
    \item \textbf{Selective Prediction:} Reject fraction $f$ of most uncertain predictions. Plot accuracy vs. $f$. Goal: steep curve (rejecting few uncertain cases improves accuracy significantly).
    
    \item \textbf{Rare Disease Performance:} Test on rare diseases (Pneumothorax: 1.1\% in IU X-Ray). Do uncertainty estimates correlate with classification errors?
\end{enumerate}

\subsubsection{Expected Contributions}

\begin{itemize}
    \item First application of conformal prediction to medical VLMs
    \item Demonstrate improved calibration (ECE reduction 20-30\%)
    \item Show 10-15\% accuracy improvement by deferring 15\% most uncertain cases
    \item Provide confidence intervals for zero-shot predictions
\end{itemize}

\textbf{Publication venue:} Radiology AI (clinical utility), MIDL (methodology)

\subsection{Extension 2: Factuality-Constrained Report Generation}

\subsubsection{Motivation}

Report generation models frequently hallucinate findings:
\begin{itemize}
    \item Generate "fracture" when none exists
    \item Omit critical findings like pneumothorax
    \item Produce anatomically impossible descriptions
\end{itemize}

BLEU/ROUGE metrics fail to detect these errors. RadGraph F1 captures factuality but is typically used only for evaluation, not during training.

\textbf{Research Question:} Can we integrate RadGraph-based factuality loss into training to reduce hallucination?

\subsubsection{Technical Approach}

\textbf{Base Model: R2Gen}

R2Gen architecture:
\begin{itemize}
    \item \textbf{Visual Encoder:} ResNet-101 (grid features) + Vision Transformer (patch features)
    \item \textbf{Memory Module:} $M$ learnable slots $\{m_1, ..., m_M\}$ storing clinical patterns
    \item \textbf{Decoder:} Transformer with cross-attention to visual features and memory
\end{itemize}

Standard training loss:
\begin{equation}
\mathcal{L}_{\text{CE}} = -\sum_{t=1}^{T} \log p(y_t | y_{<t}, x)
\end{equation}

\textbf{Proposed Extension: RadGraph-Based Factuality Loss}

\textbf{Step 1: Extract Entities During Training}

For each generated report $\hat{y}$ and reference report $y^*$:
\begin{enumerate}
    \item Apply RadGraph IE: $G_{\hat{y}} = \text{RadGraph}(\hat{y})$, $G_{y^*} = \text{RadGraph}(y^*)$
    \item $G$ = set of (entity, label) and (entity, relation, entity) tuples
\end{enumerate}

\textbf{Step 2: Define Factuality Metrics}

\textbf{Precision (Anti-Hallucination):}
\begin{equation}
P_{\text{RadGraph}} = \frac{|G_{\hat{y}} \cap G_{y^*}|}{|G_{\hat{y}}|}
\end{equation}

\textbf{Recall (Completeness):}
\begin{equation}
R_{\text{RadGraph}} = \frac{|G_{\hat{y}} \cap G_{y^*}|}{|G_{y^*}|}
\end{equation}

\textbf{F1:}
\begin{equation}
F1_{\text{RadGraph}} = \frac{2 P R}{P + R}
\end{equation}

\textbf{Step 3: Training Objective}

\textbf{Combined Loss:}
\begin{equation}
\mathcal{L}_{\text{total}} = \mathcal{L}_{\text{CE}} + \lambda (1 - F1_{\text{RadGraph}})
\end{equation}

where $\lambda$ controls factuality weight (tune on validation: $\lambda \in \{0.1, 0.5, 1.0, 2.0\}$).

\textbf{Challenge:} RadGraph extraction is non-differentiable (discrete graph operations). 

\textbf{Solution:} Use as auxiliary loss computed per-batch:
\begin{enumerate}
    \item Generate reports with beam search or sampling
    \item Extract RadGraph entities and relations
    \item Compute $F1_{\text{RadGraph}}$ on batch
    \item Use as reward in policy gradient or as soft constraint
\end{enumerate}

\textbf{Alternative (Simpler):} 
Train CheXbert classifier as auxiliary task. Add multi-task loss:
\begin{equation}
\mathcal{L}_{\text{total}} = \mathcal{L}_{\text{CE}} + \lambda \mathcal{L}_{\text{CheXbert}}
\end{equation}

where $\mathcal{L}_{\text{CheXbert}}$ is binary cross-entropy on 14 CheXpert labels extracted from generated text.

\subsubsection{Implementation Plan}

\textbf{Phase 1 (Baseline):}
\begin{enumerate}
    \item Implement R2Gen on IU X-Ray
    \item Reproduce BLEU-4 = 0.20 baseline
    \item Add RadGraph F1 evaluation (expected: 0.40-0.42)
\end{enumerate}

\textbf{Phase 2 (Extension):}
\begin{enumerate}
    \item Integrate RadGraph extraction in training loop
    \item Experiment with factuality loss variants:
    \begin{itemize}
        \item Auxiliary CheXbert loss (simpler, differentiable)
        \item Policy gradient with RadGraph F1 reward (stronger, harder to train)
        \item Retrieval augmentation (retrieve similar reports, condition generation)
    \end{itemize}
    \item Hyperparameter tuning: $\lambda$, beam size, memory slots
\end{enumerate}

\textbf{Phase 3 (Evaluation):}
\begin{enumerate}
    \item Quantitative: BLEU, RadGraph F1, CheXbert F1
    \item Qualitative: Manual analysis of 100 generated reports
    \item Error taxonomy: False positive entities, false negatives, severity errors
\end{enumerate}

\subsubsection{Expected Contributions}

\begin{itemize}
    \item Demonstrate RadGraph F1 improvement: 0.42 $\rightarrow$ 0.48 (+6\% absolute, +14\% relative)
    \item Reduce false entity rate by 20-30\%
    \item Maintain or slightly improve BLEU (factuality-fluency tradeoff)
    \item Provide error analysis: which entity types most prone to hallucination?
\end{itemize}

\textbf{Publication venue:} MICCAI, EMNLP, TMI

\subsection{Extension 3: Cross-Dataset Robustness Analysis}

\subsubsection{Motivation}

Most papers train and evaluate on single dataset. External validation rare. Unknown:
\begin{itemize}
    \item How much does performance degrade on other institutions?
    \item Which diseases generalize well vs. poorly?
    \item Are errors systematic (e.g., false positives on rare equipment artifacts)?
\end{itemize}

\textbf{Research Question:} What is the cross-dataset robustness of medical VLMs, and can simple techniques improve generalization?

\subsubsection{Experimental Protocol}

\textbf{Setup:}
\begin{enumerate}
    \item Train CheXzero on IU X-Ray (3,955 studies)
    \item Test zero-shot on:
    \begin{itemize}
        \item IU X-Ray test set (in-distribution)
        \item NIH ChestX-ray14 (out-of-distribution)
    \end{itemize}
    \item Measure AUROC, AUPRC, F1 per disease
    \item Analyze failure modes: false positives, false negatives
\end{enumerate}

\textbf{Hypotheses:}
\begin{itemize}
    \item \textbf{H1:} Performance drops 5-10\% AUROC on NIH-14 due to domain shift
    \item \textbf{H2:} Rare diseases (Pneumothorax, Pneumonia) degrade more than common diseases (Cardiomegaly, Effusion)
    \item \textbf{H3:} False positive rate increases on NIH-14 due to dataset-specific artifacts
\end{itemize}

\subsubsection{Domain Adaptation Techniques}

\textbf{Baseline: No adaptation (zero-shot transfer)}

\textbf{Technique 1: Test-Time Augmentation (TTA)}
\begin{itemize}
    \item Apply multiple augmentations (flip, rotate, crop) at test time
    \item Average predictions: $\hat{y} = \frac{1}{K} \sum_{k=1}^{K} f(x_k)$
    \item Expected improvement: +1-2\% AUROC (free, no retraining)
\end{itemize}

\textbf{Technique 2: Few-Shot Fine-Tuning}
\begin{itemize}
    \item Use 1\%, 5\%, 10\% of NIH-14 labeled data
    \item Fine-tune classification head (freeze encoder or full model)
    \item Plot: labeled data fraction vs. AUROC recovery
\end{itemize}

\textbf{Technique 3: Domain-Invariant Features}
\begin{itemize}
    \item Add domain classifier on intermediate features
    \item Adversarial training: maximize task performance, minimize domain discrimination
    \item Loss: $\mathcal{L} = \mathcal{L}_{\text{task}} - \lambda \mathcal{L}_{\text{domain}}$
\end{itemize}

\textbf{Technique 4: Ensemble}
\begin{itemize}
    \item Train multiple CheXzero models with different:
    \begin{itemize}
        \item Random seeds
        \item Data augmentation strategies
        \item Prompt templates
    \end{itemize}
    \item Aggregate predictions: majority vote or averaging
    \item Expected: +2-3\% AUROC, better calibration
\end{itemize}

\subsubsection{Evaluation}

\textbf{Quantitative Metrics:}
\begin{itemize}
    \item AUROC, AUPRC per disease
    \item Mean AUROC across diseases (macro-average)
    \item Performance drop: $\Delta = \text{AUROC}_{\text{IU}} - \text{AUROC}_{\text{NIH}}$
    \item Calibration: ECE on both datasets
\end{itemize}

\textbf{Qualitative Analysis:}
\begin{itemize}
    \item Error analysis: Sample 100 false positives and false negatives from NIH-14
    \item Identify patterns:
    \begin{itemize}
        \item Equipment artifacts (e.g., pacemaker wires misclassified as consolidation)
        \item Projection differences (AP vs. PA)
        \item Labeling errors in NIH-14 (recall: NLP-extracted labels)
    \end{itemize}
    \item Visualize: Attention maps on misclassified examples
\end{itemize}

\subsubsection{Expected Contributions}

\begin{itemize}
    \item Comprehensive cross-dataset benchmark:
    \begin{itemize}
        \item IU X-Ray: 0.85 AUROC
        \item NIH-14: 0.78 AUROC (baseline, $\Delta = 7\%$)
        \item With few-shot (5\%): 0.82 AUROC ($\Delta = 3\%$)
    \end{itemize}
    \item Identify diseases with poor generalization (rare diseases, subtle findings)
    \item Demonstrate simple techniques (TTA, few-shot) recover 50-70\% of performance drop
    \item Provide failure taxonomy for future work
\end{itemize}

\textbf{Publication venue:} MICCAI, MIDL (emphasize robustness and clinical deployment)

\subsection{Implementation Timeline}

Table \ref{tab:timeline_detailed} breaks down implementation phases:

\begin{table}[h]
\centering
\caption{Detailed Implementation Timeline}
\label{tab:timeline_detailed}
\small
\begin{tabular}{lll}
\toprule
\textbf{Phase} & \textbf{Duration} & \textbf{Tasks} \\
\midrule
\textbf{Setup} & Weeks 1-2 & Environment, datasets, reproduce baselines \\
\textbf{Extension 1} & Weeks 3-6 & Uncertainty quantification \\
\textbf{Extension 2} & Weeks 7-10 & Factuality-constrained generation \\
\textbf{Extension 3} & Weeks 11-14 & Cross-dataset evaluation \\
\textbf{Analysis} & Weeks 15-18 & Ablations, comparisons, error analysis \\
\textbf{Writing} & Weeks 19-24 & Paper draft, revision, submission \\
\bottomrule
\end{tabular}
\end{table}

\textbf{Risk Mitigation:}
\begin{itemize}
    \item If Extension 2 (factuality loss) proves too complex, focus on Extensions 1 and 3
    \item Maintain weekly progress logs and pivot if blocked
    \item Target workshop paper initially (4 pages), extend to conference (8 pages) if results strong
\end{itemize}

%=====================================================
% Note: This is first half of proposal. Second half would continue with:
% - Section 5: Experimental Plan and Evaluation
% - Section 6: Expected Contributions and Limitations
% - Section 7: Ethical Considerations
% - Section 8: Timeline and Milestones
% - Section 9: Conclusion and Future Work
% - Bibliography
%=====================================================

% Placeholder for references (will be populated from bibliography.bib)
\bibliographystyle{plain}
\bibliography{references}

\end{document}
